\documentclass[a4paper,11pt]{article}

\usepackage{fontspec}
\setmainfont{XCharter}[
  UprightFont = *-Roman,
  BoldFont = *-Bold,
  ItalicFont = *-Italic,
  BoldItalicFont = *-BoldItalic,
  Ligatures = TeX,
]
\usepackage[empty]{fullpage}
\usepackage{titlesec}
\usepackage{enumitem}
\usepackage[hidelinks]{hyperref}
\usepackage{fancyhdr}
\usepackage{xcolor}
\usepackage{multicol}

\definecolor{accent}{HTML}{1A2733}
\definecolor{subtle}{HTML}{666666}

\pagestyle{fancy}
\fancyhf{}
\renewcommand{\headrulewidth}{0pt}
\renewcommand{\footrulewidth}{0pt}

\addtolength{\oddsidemargin}{-0.5in}
\addtolength{\evensidemargin}{-0.5in}
\addtolength{\textwidth}{1.0in}
\addtolength{\topmargin}{-0.9in}
\addtolength{\textheight}{2.5in}

\urlstyle{rm}
\raggedright
\linespread{0.84}
\setlength{\parindent}{0pt}
\setlength{\tabcolsep}{0in}

\titleformat{\section}{%
  \normalfont\large\bfseries\color{accent}%
}{}{0em}{}[{\color{subtle}\vspace{1pt}\titlerule[0.4pt]}]
\titlespacing*{\section}{0pt}{6pt}{2pt}

\newcommand{\resumeSubheading}[4]{%
  \vspace{1pt}%
  \noindent{\bfseries #1}\hfill {\small #4}\par%
  \noindent{\itshape\color{subtle} #3 \,$\cdot$\, #2}\par%
}

\newcommand{\resumeProject}[3]{%
  \vspace{1pt}%
  \noindent{\bfseries #1}\hfill {\small #3}\par%
  \noindent{\itshape\color{subtle} #2}\par%
}

\setlist[description]{leftmargin=14pt, itemsep=1pt, topsep=1pt, parsep=0pt}
\setlist[itemize]{leftmargin=14pt, label=\textbullet, itemsep=0pt, topsep=0pt, parsep=0pt}

\begin{document}

%----------HEADING-----------------
\begin{center}
  {\fontsize{22pt}{26pt}\selectfont\bfseries\color{accent} Mustafa Arinmis}\\[4pt]
  {\small\color{subtle} mustafa\_arinmis@outlook.com \, $\bullet$ \, +90 543 549 99 53 \, $\bullet$ \, \href{https://github.com/arinmis}{github.com/arinmis} \, $\bullet$ \, \href{https://www.linkedin.com/in/arinmis}{linkedin.com/in/arinmis}}\\[2pt]
  {\small\color{subtle} Antalya, Türkiye \,$\cdot$\, Bilgisayar Mühendisliği Lisans Derecesi, Akdeniz Üniversitesi}
\end{center}
\vspace{1pt}%-----------SUMMARY-----------------
\section{Özet}
Backend odaklı, \textbf{4 yıllık yazılım mühendisliği deneyimine} ve uçtan uca ürün sahipliği yaklaşımına sahip \textbf{Yazılım Mühendisi}. CQRS, Dapr ve AWS EKS kullanarak 30'dan fazla servisten oluşan bir turizm platformunda \textbf{5 üretim C\#/.NET 10 mikroservisi} tasarlayıp devreye aldı. Bulut tabanlı mimari ile \textbf{RAG}, \textbf{MCP}, \textbf{LangChain} ve \textbf{Microsoft Semantic Kernel} prototipleri aracılığıyla yapay zeka entegrasyonu konusunda deneyimli.


%-----------EXPERIENCE-----------------
\section{Deneyim}

\resumeSubheading{\href{https://www.santsg.com/en/products/octopus-project/}{SAN TSG (OctopusPlus)}}{Antalya, Türkiye}{Backend Yazılım Mühendisi}{Aralık 2024 -- Temmuz 2026}
\begin{itemize}
  \item Uygulamalı turizm alanı bilgisini kullanarak \textbf{otel ve gezi rezervasyonları} için özellikler geliştirdi.
  \item Gezi, uçuş ve otel rezervasyonu alanlarında 30'dan fazla servisten oluşan turizm platformunda, \textbf{Clean/Onion mimarisi}, \textbf{CQRS}, MediatR ve \textbf{FluentValidation} kullanarak \textbf{5 adet C\#/.NET 10 üretim mikroservisini} uçtan uca geliştirip devreye aldı.
  \item Mikroservis iletişimi için \textbf{Dapr} sidecar entegrasyonunu gerçekleştirdi; \textbf{RabbitMQ} pub/sub ile \textbf{EF Core/PostgreSQL}, \textbf{MongoDB} ve durum yönetimi için \textbf{Redis} kullanan çoklu veri depolama yaklaşımını uyguladı.
  \item \textbf{30'dan fazla .NET mikroservisi} için temel bir \textbf{MCP sunucusu ve AI agent platformu} ile teknik dokümantasyona yönelik bir \textbf{RAG servisi} geliştirdi; \textbf{6 stajyere mentorluk yaptı} ve bunlardan biri yarı zamanlı mühendis olarak işe alındı.
  \item Bitbucket, JIRA ve çalışma kayıtlarından çalışan kapasite kullanım metriklerini toplayan, e-posta otomasyonu içeren ve AWS ECR'a dağıtılan otomatik bir Python servisi geliştirdi.
\end{itemize}

\resumeSubheading{Mavi ArGe}{Antalya, Türkiye}{Backend Yazılım Mühendisi (Sözleşmeli)}{Ekim 2024 -- Aralık 2024}
\begin{itemize}
  \item Bir Elektrikli Araç şarj istasyonu Filo Yönetimi \textbf{IoT} sistemi (\textbf{TÜBİTAK} projesi) için batarya telemetrisi ve izleme panolarına yönelik kural tabanlı bir motor dahil olmak üzere \textbf{C\#/.NET REST API'leri} geliştirdi.
\end{itemize}

\resumeSubheading{P.I. Works}{Uzaktan, İstanbul, Türkiye}{Full Stack Yazılım Geliştirici}{Haziran 2023 -- Eylül 2024}
\begin{itemize}
  \item Kurumsal bir telekomünikasyon projesinde \textbf{Angular-C\#/.NET} ile yeni özellikler geliştirdi, birim test kapsamını artırdı ve kod kokularını gideren yeniden düzenlemeler yaptı.
\end{itemize}

\resumeSubheading{Linarit}{Antalya, Türkiye}{Frontend Geliştirici}{Şubat 2022 -- Eylül 2022}
\begin{itemize}
  \item \textbf{React} tabanlı \textbf{Linarit PQM Endüstriyel IoT platformu} için, \textbf{5 servisli mikroservis mimarisi} ve IoT uç cihazları arasında gerçek zamanlı \textbf{gRPC} veri görselleştirmesi, istemci tarafı önbellekleme ve performans iyileştirmeleri dahil temel özellikler geliştirdi.
\end{itemize}

\resumeSubheading{AudioHealer}{Antalya, Türkiye}{Kurucu Yazılım Geliştirici}{Eylül 2021 -- Şubat 2022}
\begin{itemize}
  \item Yapay zeka destekli, full stack ses iyileştirme platformu \href{https://www.audiohealer.com}{AudioHealer}'ı geliştirdi: gerçek zamanlı gürültü giderimi için \textbf{DTLN} (TensorFlow/Keras) ve 18 bant EQ ile EBU R128 ses yüksekliği normalizasyonu için \textbf{FFmpeg} kullandı. Y Combinator'a başvurdu.
\end{itemize}

%-----------EDUCATION-----------------
\section{Eğitim}
\resumeProject{Bilgisayar Mühendisliği Lisans Derecesi}{Akdeniz Üniversitesi, Antalya}{2019 -- 2024}

%-----------PROJECTS-----------------
\section{Projeler}
\resumeProject{Tiro}{\href{https://tiro.legal}{tiro.legal}}{}
\begin{itemize}
  \item Avukatlar için yerleşik yerel LLM (RAG) içeren masaüstü dava yönetim uygulaması. \textbf{Python}, \textbf{FastAPI}, \textbf{Electron}, \textbf{React}, \textbf{SQLite}, \textbf{ChromaDB}. Platformlar arası çalıştırılabilir uygulama.
\end{itemize}

\resumeProject{FitBrand}{\href{https://fitbrand.app}{fitbrand.app}}{}
\begin{itemize}
  \item Kişisel antrenörlere özel marka kimliğiyle sunulan PWA: özel logo, renkler ve alan adıyla tamamen markalanabilir uygulamalar. Yapay zeka destekli antrenman planları, danışan yönetimi, mesajlaşma ve blog. App Store gerektirmeden \textbf{PWA} olarak yüklenir. \textbf{Next.js}, \textbf{Tamagui}, \textbf{Prisma}, \textbf{PostgreSQL}. Yapay zeka için \textbf{MCP} araçları.
\end{itemize}

\end{document}
